\hypertarget{appendix-c2a9ad1e2b49}{}
\section*{Process isolation and page permissions}
\section{Process isolation and page permissions}\label{process-isolation-and-page-permissions}

\textbf{Curriculum depth:} Introduced · \textbf{First seen:} C003 · \textbf{Latest development:} C003

\subsection{Definition and mechanism}\label{definition-and-mechanism}

An address-space mapping context determines which physical frame a virtual number names and whether read/write/execute access is permitted. Two processes may use VA 0x400000 for different frames; shared mappings can deliberately use one frame. Permission checks belong to mapping state, not to the apparent numerical validity of a pointer.

\subsection{Prerequisites}\label{prerequisites}

\begin{itemize}
\tightlist
\item
  \hyperref[appendix-18ee4f913c3c]{Addresses and virtual memory} --- Mappings give numeric locations their process-specific storage identity.
\end{itemize}

\subsection{Used by}\label{used-by}

\begin{itemize}
\tightlist
\item
  \hyperref[appendix-01e5905165c4]{Page tables and MMU} --- uses this mechanism as stated in its prerequisites.
\item
  \hyperref[appendix-b081b7abe193]{Call stack and stack frames} --- uses this mechanism as stated in its prerequisites.
\item
  \hyperref[appendix-5f9e5e2f0a5e]{Read-only data and code} --- uses this mechanism as stated in its prerequisites.
\item
  \hyperref[appendix-78cb248c8355]{ASLR} --- uses this mechanism as stated in its prerequisites.
\end{itemize}

\subsection{Recall and next depth}\label{recall-and-next-depth}

Give private and deliberately shared mapping examples.

\textbf{Pending:} Kernel/user isolation and architecture-specific protection.

\subsection{Source and navigation}\label{source-and-navigation}

\href{Continuous\%20Notes/Computing\%20Foundations\%20-\%20From\%20Bits\%20to\%20C++\%20and\%20Systems\%20APIs\#Where\%20the\%20bytes\%20live}{C003 teaching chapter} · \hyperref[appendix-fd090d982a15]{Part definitions} · \href{}{Dependency Map (02-Foundations-II, appendix-4ad47bdd269a)}
